\subsection{Nonflat even disks can have positive-form cohomology}
\label{arch136:even-disk-positive}

\paragraph{Archive and exact scope.}
This is a local algebra test for cote136, physical
p.~117/footer116/scan117 and physical p.~121/footer120/scan121.
All calculations below are explicit, tagged none. For a ring
\(A\) and an \(A\)-module \(M\), write
\[
 \mathscr D_A(M)=\Gamma_A(M[0])\otimes_A\Lambda_A(M[1]),
 \qquad du=e_u.                                         \tag{EP.1}
\]
The differential is the special PD derivation. The examples
disprove positive-form exactness for arbitrary nonflat modules.
They do not supply such a counterexample for the particular
\(\mathbb Z_{(p)}\)-free, coefficient-torsion-free bottom
lattice of the coherent realization. Moreover the displayed
classes lie in the ideal of extra closed zero-forms; this
limitation is checked explicitly at the end.

\paragraph{A cyclic arity block.}
Let \(A\) have characteristic \(p\), and let \(I\subset A\)
be an ideal with \(I^2=0\). For \(M=(A/I)u\), the
arity-\(p\) part of (EP.1) is exactly
\[
                         A\longrightarrow A/I,          \tag{EP.2}
\]
in forms zero and one, with the quotient map as differential.
Indeed the degree-\(p\) divided-power relations coming from
\(Iu=0\) vanish: a product \(iu\,\gamma_{p-1}(u)\)
has coefficient \(pi=0\), while \(\gamma_k(iu)=0\)
for \(k\geq2\). The PD ideal generated by the linear
relations has no further arity-\(p\) contribution. For
arity \(p-1\), the factorial \((p-1)!\) is invertible,
so the divided-power module is the ordinary symmetric power
\(A/I\). Tensoring with the cyclic form-one module again
gives \(A/I\). Its exterior square is zero. Finally
\(d\gamma_p(u)=\gamma_{p-1}(u)e_u\), with coefficient
one, proving (EP.2).

More generally, over \(R=\mathbb Z_{(p)}\), take a
square-zero extension \(A=R\oplus V\) and a rank-one
direct summand \(I\subset V\). Give \(V\) the PD
structure with all higher divided powers zero. The same
arity block is
\[
                         A/(pI)\longrightarrow A/I.     \tag{EP.3}
\]
Here the only arity-\(p\) relation is \(pI\), by the
same calculation; all higher powers of an element of \(I\)
are zero. Both claims can also be obtained directly from the
full cyclic presentation in ARCH-136-NONFLAT-QUOTIENT-DISK.

\paragraph{ARCH-136-EVEN-DISK-H1.}
For every prime \(p\) there is a finite PD coefficient ring
in characteristic \(p\), and a module over it, whose even
disk has nonzero cohomology in form degree one. There are also
homogeneous examples over the original coefficient ring
\(S=\Gamma_{\mathbb F_p}(T)\), with nonzero \(H^1\)
in unbounded weights for a single weightwise finite module.

\begin{proof}
First put \(A=\mathbb F_p\oplus V\), \(V^2=0\),
\(V=\mathbb F_pX\oplus\mathbb F_pY\), and let
\[
 I_1=(X),\qquad I_2=(Y),\qquad I_3=(X+Y),\qquad
                 M=\bigoplus_{i=1}^3(A/I_i)u_i.          \tag{EP.4}
\]
Give \(V\) its zero higher PD operations. The module-disk
functor sends direct sums to tensor products. In multiarity
\((p,p,p)\), (EP.2) therefore gives the tensor product of
the three complexes \(A\to A/I_i\).

Give the coefficient space \(V\) degree one solely for
this calculation. Its degree-one subcomplex is
\[
                 V\longrightarrow\bigoplus_i V/I_i
                                      \longrightarrow0. \tag{EP.5}
\]
Every two different lines span \(V\), so the next terms
have coefficient degree one equal to zero. The first map is
injective and has ranks two and three. Thus this block of
\(H^1\) is one-dimensional. In particular the class with
coordinates \((Y,0,0)\) is nonzero: a lift would have to
belong to \(I_2\cap I_3=0\) and have residue \(Y\)
modulo \(I_1\). A representative in the disk is
\[
 Z_1=Y\gamma_{p-1}(u_1)e_{u_1}
                         \gamma_p(u_2)\gamma_p(u_3).    \tag{EP.6}
\]

For the homogeneous original-coefficient version, fix
\(h=p^a\), put \(x=T^{[h]}\), \(y=T^{[ph]}\), and set
\[
 M_h=(S/(x))u\ \oplus\ (S/(y))v\ \oplus\
       (S w_0\oplus S w_1)/(xw_0+yw_1).                 \tag{EP.7}
\]
Give \(u,v,w_1\) weight \(h\) and \(w_0\) weight
\(ph\). The relation is homogeneous. The form
\[
 z_{1,h}=y\gamma_{p-1}(u)e_u
                            \gamma_p(v)\gamma_p(w_1)    \tag{EP.8}
\]
is closed. The differential of its first factor vanishes
because \(e_u^2=0\). In the term differentiating
\(\gamma_p(v)\), use \(yv=0\) and the invertibility
of \((p-1)!\) to get \(y\gamma_{p-1}(v)=0\).
In the last term, use \(yw_1=-xw_0\) and
\(x\gamma_{p-1}(u)=0\). These are actual identities,
not only equalities in cohomology.

There is a ring map from \(S\) to the finite ring in
(EP.4), sending the two indicated digit generators to
\(X,Y\) and every other digit generator to zero. Map
\(u,v\) to \(u_1,u_2\) and both \(w_0,w_1\) to
\(u_3\). The defining relations are respected. Functoriality
of the relative module divided-power and exterior algebras
gives an underlying differential algebra map carrying
(EP.8) to (EP.6). Thus (EP.8) is not a boundary.
The coefficient ring map need not preserve coefficient PD
operations, and the module map forgets weights; neither
property is needed to detect a nonboundary. The source disk
itself retains its genuine coefficient PD structure and grading.

The weight of (EP.8) is \(4ph\). Take the direct sum of
the modules (EP.7) over \(h=p^a\). Only finitely many blocks
contribute to any fixed weight, so the resulting module is
finite-dimensional in each weight. Each block is a retract
as a module, hence also as a disk. Its nonzero class survives,
and the weights are unbounded. All classes have exact
additive order \(p\).
\end{proof}

\paragraph{ARCH-136-EVEN-DISK-H2.}
The same assertions hold in form degree two. They can be
realized in multiarity \((p,p,p,p)\), and over
\(\Gamma_{\mathbb F_p}(T)\) by homogeneous six-generator
modules.

\begin{proof}
Use the square-zero ring \(A=\mathbb F_p\oplus V\), with
\(V\) having basis \(X,Y,Z\), and the four lines
\[
 I_1=(X),\ I_2=(Y),\ I_3=(Z),\ I_4=(X+Y+Z).
                                                               \tag{EP.9}
\]
Let \(M=\bigoplus_{i=1}^4(A/I_i)u_i\).
In the indicated multiarity, the coefficient-degree-one
complex has dimensions
\[
                    3\longrightarrow8\longrightarrow6
                                      \longrightarrow0. \tag{EP.10}
\]
Every three of the lines span \(V\). The first map is
injective, and its image is the kernel of the second.
To see the last assertion, first use pairwise compatibility
for the three coordinate lines to find one vector in \(V\)
having the first three given residues. Subtract it. The fourth
residue then lies in every \((I_i+I_4)/I_4\), for
\(i=1,2,3\); their intersection is zero. This proves the
assertion. The ranks of the two differentials are consequently
three and five, and the degree-two cohomology is exactly
one-dimensional.

Here is a detector for a specific class, including the signs.
Use the ordered pair coordinates
\((12),(13),(23),(14),(24),(34)\), represented respectively
by \(Z,Y,X,Y,X,X\) in their one-dimensional quotients.
The functional
\[
                (12)-(13)+(23)+(14)-(24)+(34)            \tag{EP.11}
\]
vanishes on the image of the degree-one differential
\((v_i)_i\mapsto(v_j-v_i)_{i<j}\). This follows by
substitution of the two coordinates of each residue, or by
canceling each of them separately. It takes value one on
the vector with only coordinate \((12)=Z\). Thus
\[
 Z_2=Z\gamma_{p-1}(u_1)e_{u_1}
       \gamma_{p-1}(u_2)e_{u_2}\gamma_p(u_3)\gamma_p(u_4)
                                                               \tag{EP.12}
\]
represents the nonzero class. Its differential vanishes
using \(Zu_3=0\) and
\(Zu_4=-(X+Y)u_4\); the two resulting terms are killed
by \(Xu_1=0\) and \(Yu_2=0\).

For the homogeneous version, put
\(x=T^{[h]}\), \(y=T^{[ph]}\), \(z=T^{[p^2h]}\),
where \(h=p^a\), and set
\[
 \begin{split}
 M_h={}&(S/(x))u\oplus(S/(y))v\oplus(S/(z))t\\
 &\oplus(Sw_0\oplus Sw_1\oplus Sw_2)/
                              (xw_0+yw_1+zw_2).
 \end{split}                                             \tag{EP.13}
\]
Give \(u,v,t,w_2\) weight \(h\), give \(w_0\) weight
\(p^2h\), and give \(w_1\) weight \((p^2-p+1)h\).
The relation is homogeneous. The closed form is
\[
 z_{2,h}=z\gamma_{p-1}(u)e_u\gamma_{p-1}(v)e_v
                                  \gamma_p(t)\gamma_p(w_2).
                                                               \tag{EP.14}
\]
The same special-derivation calculation applies: differentiating
\(\gamma_p(t)\) is killed by \(zt=0\), and
\(zw_2=-xw_0-yw_1\) kills the last term using the two
first factors. Map the three digit coefficients to \(X,Y,Z\)
in the finite ring, and map \(w_0,w_1,w_2\) all to
\(u_4\). The underlying module-disk map sends (EP.14)
to (EP.12), proving nonvanishing. Its weight is
\((p^2+4p)h\). Direct sums over \(h=p^a\) give a single
weightwise finite module and unbounded nonzero order-\(p\)
classes, as before.
\end{proof}

\paragraph{An integral finite-coefficient variant, and its limitation.}
Both finite examples also work with
\(R=\mathbb Z_{(p)}\) and \(A=R\oplus V\), \(V^2=0\),
using exactly the displayed primitive rank-one submodules
\(I_i\subset V\). All the modules \(A/I_i\), and hence
their direct sums, are free over \(R\). Formula (EP.3)
applies. In coefficient degree one, tensoring the arity-\(p\)
blocks gives
\[
 V\big/\left(\sum_{i\in J} I_i+
                       p\sum_{i\notin J}I_i\right)
                   \quad\text{in the component }J.      \tag{EP.15}
\]
Since all the lines together span \(V\), this is
\((V/pV)/\sum_{i\in J}\overline I_i\). Thus the entire
coefficient-degree-one complex is exactly its previously
computed \(\mathbb F_p\) complex. The same representatives
give order-\(p\) \(H^1\) and \(H^2\), respectively.
This supplies genuinely \(R\)-free nonflat parameter modules,
but over a different coefficient ring with square-zero elements.
It is not an example over the domain \(\Gamma_R(T)\), and
does not claim the saturated-lattice property (AT.2).

\paragraph{Why these classes do not yet obstruct low-degree repair.}
The displayed positive classes are products of extra closed
zero-forms with other forms. In (EP.6), the factor
\(Y\gamma_p(u_2)\) is closed because \(Yu_2=0\) and
\((p-1)!\) is invertible in the characteristic-\(p\) example.
In (EP.12), use the factor \(Z\gamma_p(u_3)\).
The analogous factors \(y\gamma_p(v)\) and
\(z\gamma_p(t)\) work in the homogeneous original-coefficient
examples. The integral finite-coefficient variants have the same
closed factors, since the linear annihilator identities imply
the needed derivative vanishings before reducing modulo \(p\).
All these zero-forms have positive parameter arity.

Consequently the positive classes exhibited here are killed
by quotienting out the ideal generated by the extra closed
zero-forms. The calculation disproves a positive-form
Poincare lemma for the raw nonflat disk, but neither proves
nor disproves exactness of that subsequent quotient. Nor does
it determine what new classes the quotient or the full
AQ-136-RELATIVE-H1 iteration might introduce. In particular
the \(\mathbb Z_{(p)}\)-free, coefficient-torsion-free
actual bottom lattice, after its low-degree repair, remains
outside the claimed counterexamples.

\paragraph{Verification status.}
The author has read the complete saved module, including the cyclic
blocks, both finite complexes, all-prime homogeneous detectors,
integral finite-coefficient variants, and the closed-zeroform
factorizations. The parent and reviewers123 and1348 have each read
the complete saved proof and communicated PASS on 2026-11-22. The finite program
\texttt{even\_disk\_probe.py} was used only for exploration
at \(p=2\); no statement above relies on a finite computation.
